\documentclass[5p,times]{elsarticle}
\usepackage{amssymb}
\usepackage{amsmath}
\journal{Journal of Fictional Ecology}
\begin{document}
\begin{frontmatter}
\title{Pollinator networks along an urbanisation gradient in a fictional river city}
\author[ku]{Sakura Minase\corref{cor1}}
\ead{minase@example.org}
\author[nw]{Oliver Penhallow}
\author[ku]{Ren Takamura}
\cortext[cor1]{Corresponding author}
\affiliation[ku]{organization={Graduate School of Environmental Science, Kitamachi University},
            city={Kitamachi},
            postcode={100-0000},
            country={Japan}}
\affiliation[nw]{organization={School of Biological Fictions, University of Nowhere},
            city={Nowhere},
            postcode={00000},
            country={Freedonia}}
\begin{abstract}
Urbanisation transforms the habitats available to pollinators, but its effect on the structure of plant--pollinator networks remains debated.
We sampled interactions between flowering plants and insect visitors at 24 sites along a fictional river that runs from farmland into a dense city centre.
Species richness of pollinators declined by about 40\% towards the centre, whereas the richness of flowering plants remained stable because of ornamental planting.
Networks in the city centre were smaller and more nested, and they relied on a few generalist bee species.
Simulated species loss showed that these urban networks are robust to random extinctions but fragile when the most connected pollinators disappear.
Our results suggest that conserving a small set of generalist pollinators is key to maintaining pollination services in cities.
\end{abstract}
\begin{keyword}
pollination \sep urban ecology \sep ecological networks \sep nestedness
\end{keyword}
\end{frontmatter}
\section{Introduction}
Cities now cover a growing share of the land surface, and their effects on biodiversity are a central question in ecology \cite{minase2021,penhallow2022}.
Pollinators are of particular concern because they provide an ecosystem service on which many crops and wild plants depend.
Studies of individual species show mixed responses to urbanisation, but the consequences for whole interaction networks are less clear \cite{takamura2023}.

Here we use a fictional river corridor as a natural gradient from rural to urban land use.
We ask how the size and structure of plant--pollinator networks change along this gradient and how robust the networks are to species loss.
We expected smaller and less specialised networks in the city centre, but we also expected ornamental plants to buffer some of these effects.

\section{Methods}
\subsection{Study sites}
We selected 24 sites along the river.
\bibliographystyle{elsarticle-num}
\bibliography{refs}
\end{document}
